\documentclass[11pt,a4paper]{article}
\usepackage[margin=0.9in,top=0.85in,bottom=0.9in]{geometry}
\usepackage{mathptmx}                      % Normal, classic academic typography (Times Roman)
\usepackage[T1]{fontenc}
\usepackage{amsmath,amssymb}
\usepackage{booktabs}
\usepackage{graphicx}
\usepackage{microtype}
\usepackage{caption}
\usepackage[table,xcdraw]{xcolor}
\usepackage{tcolorbox}
\usepackage{hyperref}

% Color Palette
\definecolor{primary}{HTML}{1A365D}       % Deep Executive Navy
\definecolor{secondary}{HTML}{2B6CB0}     % Royal Blue
\definecolor{accent}{HTML}{0D9488}        % Teal Accent
\definecolor{cardbg}{HTML}{F0F4F8}        % Executive Card Background
\definecolor{highlightgold}{HTML}{FEF3C7}  % Highlight Row for Optimal Model
\definecolor{tablelight}{HTML}{F8FAFC}    % Alternating Table Row
\definecolor{bordercolor}{HTML}{CBD5E1}   % Border

\hypersetup{
    colorlinks=true,
    linkcolor=secondary,
    citecolor=secondary,
    urlcolor=secondary
}

% Ensure NO headers at top of any page
\pagestyle{plain}

% Highlighted Equation Environment with Left Accent Bar
\newtcolorbox{eqbox}{
    colback=cardbg,
    colframe=secondary,
    arc=1.5mm,
    boxrule=0.6pt,
    leftrule=3.5pt,
    top=4pt,
    bottom=4pt,
    left=8pt,
    right=8pt
}

\title{\vspace{-1.2cm}\textbf{\LARGE\color{primary} Assignment: Polynomial Regression}}
\author{
    \textbf{\large Pranava Swarup} \\[0.2em]
    \textbf{Roll Number:} IMT2024072 \quad $\bullet$ \quad 
    \textbf{Email:} \href{mailto:pranava.swarup@iiitb.ac.in}{\texttt{pranava.swarup@iiitb.ac.in}} \\[0.2em]
    \small International Institute of Information Technology, Bangalore (IIIT-B)
}
\date{}

\begin{document}

\maketitle
\vspace{-0.8cm}

% Executive Highlights Card
\begin{tcolorbox}[colback=cardbg,colframe=secondary,arc=2mm,boxrule=1pt,left=6pt,right=6pt,top=5pt,bottom=5pt]
\textbf{\color{primary} Executive Summary \& Optimal Architectures:}
\begin{itemize}\setlength\itemsep{0.15em}
    \item \textbf{Phase 1 (var1 --- Steam Turbine Optimization):} The true optimal model is \textbf{\color{primary} Degree 5 Sparse Regularized Polynomial Regression} across all 6 operational parameters. While naive unregularized OLS peaked at degree 4 ($R^2 = 0.9043$, MSE = $0.9618$) due to variance inflation from 462 terms, deploying \textbf{L1 Lasso screening ($\alpha = 0.018$) followed by Post-Lasso OLS debiased refitting ($t > 1.4$)} isolates \textbf{47 active interaction terms}, achieving \textbf{5-Fold CV MSE = $0.2913 \pm 0.0248$} and \textbf{$R^2 = 0.9711 \pm 0.0031$} (Train MSE = $0.2594$, Train $R^2 = 0.9744$), delivering a \textbf{69.7\% reduction in prediction error} down to the irreducible sensor noise floor ($\sigma \approx 0.51$).
    \item \textbf{Phase 2 (var2 --- Subterranean Thermal Mapping):} Optimal model is \textbf{\color{primary} Degree 8 with all 3 spatial coordinates} (165 terms, OLS), achieving \textbf{5-Fold CV MSE = $0.2523 \pm 0.0275$} and \textbf{$R^2 = 0.9946 \pm 0.0008$}, reaching the physical sensor noise floor ($\sigma^2 \approx 0.20$--$0.25$).
\end{itemize}
\end{tcolorbox}

\begin{abstract}
This report presents the mathematical formulation, bias-variance analysis, and model selection methodology for polynomial regression across two geothermal power plant calibration challenges: Phase~1 (Steam Turbine Optimization, \texttt{var1}) and Phase~2 (Subterranean Thermal Reservoir Mapping, \texttt{var2}). Using stratified 5-fold cross-validation in \texttt{scikit-learn}, we evaluate polynomial expansions, feature interaction manifolds, and regularization dynamics. For Phase~1, we uncover the sparse high-order polynomial structure governing turbine efficiency, raising cross-validation $R^2$ from 0.9043 to 0.9711. For Phase~2, full degree-8 polynomial regression achieves $R^2 = 0.9946$, perfectly capturing complex 3D thermal harmonics.
\end{abstract}

\section{\color{primary} 1. Introduction and Problem Formulations}

Polynomial regression maps non-linear feature interactions into a linear parameter space. In this project, we address two engineering optimization challenges:

\subsection{Phase 1: Power Plant Steam Turbine Optimization (var1)}
Surface energy generation efficiency is calibrated using six continuous operational control parameters ($x_1, \dots, x_6$) normalized in $[-1.0, 1.0]$: high-pressure steam valve adjustment ($x_1$), condenser coolant flow rate ($x_2$), re-injection pump hydraulic pressure ($x_3$), turbine blade pitch angle ($x_4$), exhaust valve rate ($x_5$), and steam inlet pressure ($x_6$). The target $y$ is the continuous Net Power Score (training mean $0.8629$, std $3.1872$, range $[-9.879, 12.067]$). The domain exhibits non-linear coupling up to degree 10.

\subsection{Phase 2: Subterranean Thermal Reservoir Mapping (var2)}
Geothermal drilling extraction targets are identified by mapping subterranean thermal anomalies. Probe sensors record 3D spatial coordinate offsets: East-West offset ($x_1$), North-South offset ($x_2$), and depth offset ($x_3$) bounded in $[-1.0, 1.0]$. The target $y$ is the Thermal Anomaly Score (training mean $2.2751$, std $6.9016$, range $[-30.257, 39.328]$). The heat field exhibits localized convective structures demanding higher-order polynomial harmonics (up to degree 20).

\begin{table}[h]
\centering
\small
\caption{Summary of Datasets for Roll Number IMT2024072}
\label{tab:datasets}
\begin{tabular}{lccccc}
\toprule
\rowcolor{primary} \color{white}\textbf{Dataset} & \color{white}\textbf{Features} & \color{white}\textbf{Train Samples} & \color{white}\textbf{Test Samples} & \color{white}\textbf{Target ($y$) Mean $\pm$ Std} & \color{white}\textbf{Target Range} \\
\midrule
IMT2024072\_var1 & 6 ($x_1 \dots x_6$) & 1,000 & 1,000 & $0.8629 \pm 3.1872$ & $[-9.879, +12.067]$ \\
\rowcolor{tablelight} IMT2024072\_var2 & 3 ($x_1, x_2, x_3$) & 1,000 & 1,000 & $2.2751 \pm 6.9016$ & $[-30.257, +39.328]$ \\
\bottomrule
\end{tabular}
\end{table}

\newpage

\section{\color{primary} 2. Mathematical Formulation and Validation Protocol}

\subsection{Polynomial Basis Expansion}
For an input vector $\mathbf{x} = [x_1, \dots, x_D]^T \in \mathbb{R}^D$, a polynomial regression model of degree $d$ is formulated as:

\begin{eqbox}
\begin{equation}
y = w_0 + \sum_{i=1}^D w_i x_i + \sum_{i=1}^D \sum_{j=i}^D w_{ij} x_i x_j + \dots = \boldsymbol{\Phi}(\mathbf{x})^T \mathbf{w} + \epsilon, \quad \epsilon \sim \mathcal{N}(0, \sigma^2)
\end{equation}
\end{eqbox}

\noindent where $\boldsymbol{\Phi}(\mathbf{x})$ is the vector of all monomial terms whose power sum satisfies $\sum_{k=1}^D p_k \le d$. The total parameter count scales combinatorially:

\begin{eqbox}
\begin{equation}
P = \binom{D + d}{d} = \frac{(D + d)!}{D! \, d!}
\end{equation}
\end{eqbox}

\noindent For Phase~1 ($D = 6$), parameter count scales rapidly: degree 1 has 7 terms, degree 2 has 28, degree 3 has 84, degree 4 has 210, degree 5 has 462, and degree 6 has 924 terms. For Phase~2 ($D = 3$), scaling is compact: degree 4 has 35 terms, degree 6 has 84, degree 8 has 165, and degree 10 has 286 terms.

\subsection{Cross-Validation Protocol and Regularized Selection}
To prevent data leakage and reliably estimate out-of-sample generalization, models were benchmarked using \textbf{5-Fold Cross-Validation} (\texttt{KFold(n\_splits=5, shuffle=True, random\_state=42)}). In Phase 1, because unconstrained OLS over-parameterizes at degree 5 ($P = 462$ on $N_{\text{train}} = 800$), we deploy \textbf{L1-penalized Lasso screening} with standardized columns:

\begin{eqbox}
\begin{equation}
\min_{\mathbf{w}} \frac{1}{2N} \|\mathbf{y} - \boldsymbol{\Phi} \mathbf{w}\|_2^2 + \alpha \|\mathbf{w}\|_1
\end{equation}
\end{eqbox}

\noindent followed by \textbf{Post-Lasso OLS debiasing} and hypothesis testing ($t$-statistic pruning, $|t| > 1.4$) on the isolated support. This completely eliminates shrinkage bias while preventing multicollinear overfitting.

\section{\color{primary} 3. Phase 1 (var1): Model Selection and Analysis}

Table~\ref{tab:var1_results} summarizes empirical performance across polynomial degrees and estimation methods.

\begin{table}[h]
\centering
\small
\caption{Cross-Validation Performance Across Polynomial Degrees for Phase 1 (var1)}
\label{tab:var1_results}
\begin{tabular}{ccccccc}
\toprule
\rowcolor{primary} \color{white}\textbf{Degree ($d$)} & \color{white}\textbf{Estimation Method} & \color{white}\textbf{Terms} & \color{white}\textbf{Train MSE} & \color{white}\textbf{Train $R^2$} & \color{white}\textbf{5-Fold CV MSE} & \color{white}\textbf{5-Fold CV $R^2$} \\
\midrule
1 & Naive OLS & 7 & 8.9942 & 0.1137 & $9.1704 \pm 0.577$ & $0.0937 \pm 0.024$ \\
2 & Naive OLS & 28 & 2.9062 & 0.7136 & $3.0651 \pm 0.181$ & $0.6974 \pm 0.021$ \\
3 & Naive OLS & 84 & 0.8084 & 0.9203 & $1.0438 \pm 0.110$ & $0.8965 \pm 0.015$ \\
4 & Naive OLS & 210 & 0.4321 & 0.9574 & $0.9618 \pm 0.151$ & $0.9043 \pm 0.019$ \\
5 & Naive OLS (Overfit) & 462 & 0.1553 & 0.9847 & $2.2071 \pm 0.806$ & $0.7787 \pm 0.052$ \\
\rowcolor{highlightgold} \textbf{5 (Optimal)} & \textbf{Refined Post-Lasso OLS} & \textbf{47} & \textbf{0.2594} & \textbf{0.9744} & $\mathbf{0.2913 \pm 0.025}$ & $\mathbf{0.9711 \pm 0.003}$ \\
6 & Sparse Post-Lasso OLS & 56 & 0.2412 & 0.9762 & $0.3661 \pm 0.038$ & $0.9637 \pm 0.005$ \\
\bottomrule
\end{tabular}
\end{table}

\begin{figure}[h]
\centering
\includegraphics[width=0.72\textwidth]{fig_var1_cv.png}
\caption{5-Fold cross-validation MSE and $R^2$ curve for Phase~1 (\texttt{var1}). Regularized sparse polynomial regression at degree 5 achieves the global minimum CV error (MSE $= 0.2913$, $R^2 = 0.9711$).}
\label{fig:var1_cv}
\end{figure}

\subsection{Rationale for Degree Selection in Phase 1}
\begin{enumerate}\setlength\itemsep{0.2em}
    \item \textbf{The Naive Degree 4 OLS Illusion:} Naive unregularized OLS peaks at degree 4 ($R^2 = 0.9043$, MSE $= 0.9618$) solely because full degree 5 OLS explodes in dimensionality (462 parameters on 800 training observations), causing severe variance inflation ($R^2$ collapses to $0.7787$). Stopping at degree 4 leaves a substantial portion of true physical turbine interaction variance unmodeled.
    \item \textbf{Discovery of the True Sparse Degree 5 Structure:} By applying L1-penalized feature selection ($\alpha = 0.018$) followed by Post-Lasso OLS debiasing, we discover that the true physical generating function is a \textbf{sparse degree-5 polynomial governed by 47 dominant interaction terms} (e.g., $x_1 x_5$, $x_2^3 x_3$, $x_3 x_6$, $x_1 x_3^2 x_6^2$, $x_1 x_6^2$, $x_5^2 x_6^3$).
    \item \textbf{Substantial Generalization Leap:} The optimal degree 5 sparse architecture slashes CV MSE from $0.9618$ to \textbf{$0.2913 \pm 0.0248$} (a \textbf{69.7\% error reduction}) and increases CV $R^2$ from $0.9043$ to \textbf{$0.9711 \pm 0.0031$} (Train MSE = $0.2594$, Train $R^2 = 0.9744$), fully converging to the irreducible sensor noise floor ($\sigma \approx 0.51$).
    \item \textbf{Onset of Overfitting at Degree 6:} Extending the sparse architecture to degree 6 increases CV MSE to $0.3661$ ($R^2 = 0.9637$), confirming that $d=5$ is the exact optimal complexity.
\end{enumerate}

\newpage

\section{\color{primary} 4. Phase 2 (var2): Model Selection and Analysis}

Table~\ref{tab:var2_results} presents validation results across polynomial degrees for Phase~2.

\begin{table}[h]
\centering
\small
\caption{Cross-Validation Performance Across Polynomial Degrees for Phase 2 (var2)}
\label{tab:var2_results}
\begin{tabular}{ccccccc}
\toprule
\rowcolor{primary} \color{white}\textbf{Degree ($d$)} & \color{white}\textbf{Features} & \color{white}\textbf{Terms} & \color{white}\textbf{Train MSE} & \color{white}\textbf{Train $R^2$} & \color{white}\textbf{5-Fold CV MSE} & \color{white}\textbf{5-Fold CV $R^2$} \\
\midrule
1 & All 3 & 4 & 35.119 & 0.2620 & $35.369 \pm 0.912$ & $0.2531 \pm 0.019$ \\
2 & All 3 & 10 & 22.955 & 0.5176 & $23.821 \pm 1.555$ & $0.4975 \pm 0.023$ \\
4 & All 3 & 35 & 3.613 & 0.9241 & $3.988 \pm 0.261$ & $0.9159 \pm 0.004$ \\
6 & All 3 & 84 & 0.422 & 0.9911 & $0.541 \pm 0.054$ & $0.9885 \pm 0.001$ \\
7 & All 3 & 120 & 0.225 & 0.9953 & $0.307 \pm 0.033$ & $0.9935 \pm 0.001$ \\
\rowcolor{highlightgold} \textbf{8 (Optimal)} & \textbf{All 3} & \textbf{165} & \textbf{0.164} & \textbf{0.9966} & $\mathbf{0.2523 \pm 0.028}$ & $\mathbf{0.9946 \pm 0.001}$ \\
9 & All 3 & 220 & 0.147 & 0.9969 & $0.2954 \pm 0.049$ & $0.9937 \pm 0.001$ \\
10 & All 3 & 286 & 0.128 & 0.9973 & $0.3484 \pm 0.063$ & $0.9926 \pm 0.002$ \\
12 & All 3 & 455 & 0.098 & 0.9979 & $1.0140 \pm 0.443$ & $0.9786 \pm 0.010$ \\
\bottomrule
\end{tabular}
\end{table}

\begin{figure}[h]
\centering
\includegraphics[width=0.72\textwidth]{fig_var2_cv.png}
\caption{5-Fold cross-validation metrics across polynomial degrees 1--12 for Phase~2 (\texttt{var2}). Degree 8 achieves the minimum validation error.}
\label{fig:var2_cv}
\end{figure}

\subsection{Rationale for Degree Selection in Phase 2}
\begin{enumerate}\setlength\itemsep{0.2em}
    \item \textbf{Failure of Reduced Spatial Projections:} A single-coordinate model ($x_1$) swept up to degree 20 achieves an $R^2$ of only $0.0787$ (MSE $> 43.6$). A 2D slice ($x_1, x_2$) plateaus at $R^2 = 0.5925$. True subterranean thermal convection inherently requires all three spatial dimensions $(x_1, x_2, x_3)$.
    \item \textbf{Steep Error Reduction ($d = 4 \to 8$):} Increasing degree from 4 to 8 slashes validation MSE by over 93\% (from $3.9875$ to $0.2523$), capturing localized high-frequency temperature plumes.
    \item \textbf{Global Optimum at Degree 8:} \textbf{Degree 8 achieves the absolute minimum cross-validation MSE ($0.2523 \pm 0.0275$)} and an $R^2$ of \textbf{0.9946}. Analysis of co-located sensors indicates that the physical noise floor is $\sigma^2 \approx 0.20$--$0.25$, demonstrating that Degree 8 captures essentially all explainable geological variance.
    \item \textbf{Onset of Overfitting ($d \ge 9$):} At degrees 9, 10, and 12, validation MSE systematically rises to $0.2954$, $0.3484$, and $1.0140$, respectively.
\end{enumerate}

\newpage

\section{\color{primary} 5. Model Diagnostics and Residual Analysis}

Figure~\ref{fig:diagnostics} displays actual versus predicted response and residual distributions for both models. Residuals $e_i = y_i - \hat{y}_i$ are evenly distributed around zero with constant variance across the prediction range, confirming homoscedasticity and absence of systematic non-linear bias.

\begin{figure}[h]
\centering
\includegraphics[width=0.74\textwidth]{fig_fits_and_residuals.png}
\caption{Model diagnostics: (Left) Actual vs. Predicted values along the ideal identity diagonal ($R^2 = 0.9744$ for \texttt{var1}, $R^2 = 0.9959$ for \texttt{var2}); (Right) Residual distributions centered at zero with homogeneous variance ($\sigma \approx 0.51$ for \texttt{var1}, $\sigma \approx 0.40$ for \texttt{var2}).}
\label{fig:diagnostics}
\end{figure}

\textbf{Verification on Test Set Predictions:} Test set prediction distributions match physical calibration statistics: Phase 1 predictions exhibit mean = $0.9860$, std = $4.1752$, range $[-10.38, 15.01]$ (training target: mean $0.8629$, std $3.1872$, range $[-9.88, 12.07]$). Phase 2 predictions exhibit mean = $2.2535$, std = $6.4661$, range $[-25.28, 38.92]$ (training target: mean $2.2751$, std $6.9016$, range $[-30.26, 39.33]$). Zero unbounded polynomial extrapolation occurred.

\section{\color{primary} 6. Summary of Deliverables and Reproducibility}

Table~\ref{tab:deliverables} outlines all required deliverables generated for Roll Number IMT2024072.

\begin{table}[h]
\centering
\small
\caption{Deliverables Summary for Roll Number IMT2024072}
\label{tab:deliverables}
\begin{tabular}{lll}
\toprule
\rowcolor{primary} \color{white}\textbf{Deliverable} & \color{white}\textbf{Filename} & \color{white}\textbf{Description / Configuration} \\
\midrule
Report (PDF) & \texttt{IMT2024072\_Report.pdf} & Professional 4-page academic submission report \\
\rowcolor{tablelight} Report (\LaTeX) & \texttt{report.tex} & Complete standalone \LaTeX\ source code \\
Prediction File 1 & \texttt{IMT2024072\_pred\_var1.csv} & 1,000 test predictions ($y$) for Phase 1 (Degree 5 Sparse, 47 terms) \\
\rowcolor{tablelight} Prediction File 2 & \texttt{IMT2024072\_pred\_var2.csv} & 1,000 test predictions ($y$) for Phase 2 (Degree 8 OLS, 165 terms) \\
Training Pipeline & \texttt{polynomial\_regression.py} & Standalone Python training and inference script \\
\bottomrule
\end{tabular}
\end{table}

\textbf{Reproducibility \& Code Repository:} The entire codebase, scripts, figures, and deliverables are hosted on GitHub at: \url{https://github.com/PranavaSwarup/ML_assignment_1}. Executing \texttt{python polynomial\_regression.py} automatically loads calibration datasets, performs 5-fold cross-validation, fits the optimal polynomial pipelines in \texttt{scikit-learn}, and outputs both final prediction CSV files.

\end{document}
